% Осень26!B22; нумерация внутри столбца без пустых строк.
% Источник: лекция 4, §3 (строка 91)
\definitionticket{19}{Борелевская сигма-алгебра}
\label{definition:19}

\subsection*{Борелевская σ-алгебра}

Мотивировка: мера была построена сначала на полукольце ячеек, затем продолжена на σ-алгебру. Открытое множество --- «более универсальный объект, чем ячейка»: топологическое понятие, работающее всюду.

\begin{definition}

\emph{борелевская σ-алгебра.} Пусть \(X\) --- топологическое пространство (можно думать: метрическое; например \(\mathbb{R}^m\) с открытыми множествами, определёнными через метрику). \textbf{Борелевская σ-алгебра} \(\mathcal{B}(X)\) --- наименьшая по включению σ-алгебра, содержащая все открытые множества: \[\mathcal{B}(X):=\bigcap\{\mathcal{C}\subset 2^X : \mathcal{C}\ \text{— σ-алгебра},\ \mathcal{C}\supset\{\text{все открытые множества}\}\}.\] Корректность: пересечение любого семейства σ-алгебр --- снова σ-алгебра, поэтому минимальная σ-алгебра существует.

\end{definition}

\textbf{Замечания.}

\begin{itemize}
\tightlist
\item
  \(\mathcal{B}(X)\) содержит и все замкнутые множества (дополнения открытых; операция законна). Внимание обычно концентрируют на открытых.
\item
  На практике конструкция означает: счётные пересечения открытых, объединения таких, снова пересечения --- «в общем, довольно хитрый объект»; лучше думать о формальном определении через пересечение всех σ-алгебр.
\item
  \emph{Спойлер (обоснование позже):} лебеговская σ-алгебра \(\mathcal{M}^m\) \textbf{шире} борелевской: \(\mathcal{B}(\mathbb{R}^m)\subsetneq\mathcal{M}^m\). Борелевская σ-алгебра «плохо дружит» с полнотой: на ней нельзя определить полную меру, а мера Лебега по теореме о продолжении \textbf{полна} --- вместе с любым множеством меры ноль измеримы все его подмножества. Из-за этой «несуразицы» борелевской σ-алгебры «маловато для содержательной теории меры».
\end{itemize}
